\section{The Euler product as an analytic determinant}
\label{sec:euler-det}

There are two Hilbert spaces in an Euler product: one has a vector for
each prime, and the other has a vector for each occupation configuration.
For the two primes \(2,3\), put
\[
 K_{\{2,3\}}(s)=\operatorname{diag}(2^{-s},3^{-s})
 \quad\text{on }\mathbb C^2.
\]
The occupation basis is \(e_{a,b}\), with \(a,b\in\mathbb N_0\).
Its energy is \(a\log2+b\log3\), so for \(\operatorname{Re}s>0\),
\[
 \sum_{a,b\geq0}2^{-as}3^{-bs}
 =\det(I-K_{\{2,3\}}(s))^{-1}.
\]
At \(s=1\), the trace is \(3\) and the determinant is \(1/3\).
Passing to all primes requires convergence in an operator ideal.

\subsection{The prime operator and its Fock space}

Let \(\mathscr K=\ell^2(\mathcal P)\), where \(\mathcal P\) is the set
of rational primes.  Define
\[
 Ae_p=(\log p)e_p,\qquad
 \mathcal D(A)=\left\{x\in\mathscr K:
 \sum_p(\log p)^2|x_p|^2<\infty\right\}.
\]
The inner product is linear in its first argument.  This diagonal
operator is closed: convergence of a sequence and of its images gives
coordinatewise convergence, and the limiting image has coordinates
\((\log p)x_p\), placing the limit in the displayed domain.
Its adjoint has the same domain, since testing against the finitely
supported vectors forces those coordinates and their square summability.
Thus \(A\) is positive and self-adjoint.  The resolvent
\((A+I)^{-1}\) is the norm limit of finite diagonals because its
entries \(1/(1+\log p)\) tend to zero; hence it is compact.
For \(\operatorname{Re}s>0\), let
\[
 K(s)=e^{-sA},\qquad K(s)e_p=p^{-s}e_p.
\]
For an integer \(m\geq1\), the Schatten ideal \(\mathcal S_m\) consists
of compact operators with \(m\)-summable singular values.  The singular
values \(\mu_j(T)\) are the nonnegative square roots of the eigenvalues
of \(T^*T\), counted with multiplicity.  The ideal norm is
\(\|T\|_m=(\sum_j\mu_j(T)^m)^{1/m}\); \(\mathcal S_1\) is the
trace class.  For a trace-class diagonal \(Te_j=t_je_j\), the
Fredholm determinant used here is the convergent product
\(\det(I-T)=\prod_j(1-t_j)\).  Absolute summability of the entries
makes this the limit of the finite diagonal determinants.

\begin{theorem}[Analytic Euler determinant]
\label{thm:euler-fredholm}
For \(m\geq1\), the map \(s\mapsto K(s)\) is holomorphic with values
in \(\mathcal S_m\) precisely on the half-plane
\(\operatorname{Re}s>1/m\).  There,
\[
 \|K(s)\|_m^m=\sum_p p^{-m\operatorname{Re}s}.
\]
Let
\(\mathscr F_{\mathrm{sym}}(\mathscr K)=\bigoplus_{q\geq0}\operatorname{Sym}^q
\mathscr K\) be the Hilbert bosonic Fock space, with normalized occupation
basis \(|(r_p)_p\rangle\) indexed by finitely supported nonnegative
integers.  The summand \(q\) has particle number \(q=\sum_p r_p\);
the vacuum is the unit vector in the \(q=0\) summand.  The map
\[
 U:\mathscr F_{\mathrm{sym}}(\mathscr K)\longrightarrow\ell^2(\mathbb N^\times),
 \qquad U|(r_p)_p\rangle=e_{\prod_p p^{r_p}}
\]
is unitary.  Second quantization satisfies
\[
 U\left(\bigoplus_{q\geq0}K(s)^{\otimes_s q}\right)U^{-1}
 =e^{-sH},\qquad \operatorname{Re}s>0.
\]
These Fock operators are trace class exactly when
\(\operatorname{Re}s>1\), and on that half-plane
\[
 \operatorname{Tr}_{\mathscr F_{\mathrm{sym}}(\mathscr K)}
 \left(\bigoplus_{q\geq0}K(s)^{\otimes_s q}\right)
 =\det(I-K(s))^{-1}=\zeta(s).
\]
The determinant is the ordinary Fredholm determinant on \(\mathscr K\),
normalized to tend to \(1\) as real \(s\to+\infty\).
\end{theorem}

\begin{proof}
The singular values of \(K(s)\) are \(p^{-\operatorname{Re}s}\).
\label{proof:prime-harmonic}
The prime series \(\sum_p p^{-x}\) converges for \(x>1\) by comparison
with \(\sum_{n\geq2}n^{-x}\).  It diverges for \(x\leq1\).
Indeed, convergence of \(\sum_p1/p\) would bound all products
\(\prod_{p\leq R}(1-1/p)^{-1}\): the series
\(-\log(1-1/p)\) is at most \(2/p\).  Their geometric expansions contain
\(\sum_{n\leq R}1/n\), which is unbounded.  Comparison proves the
remaining divergence claims.

For a compact subset of \(\operatorname{Re}s>1/m\), choose
\(\sigma_0>1/m\) below its real parts.  The finite diagonal truncations
of \(K(s)\) converge uniformly there in \(\mathcal S_m\), because their
\(m\)-norm tails are bounded by \(\sum_{p>R}p^{-m\sigma_0}\).
They are holomorphic, so their limit is holomorphic in that Banach space.
At \(\operatorname{Re}s=0\), the diagonal operator is unitary and
noncompact.  For \(\operatorname{Re}s<0\), it is unbounded on its
maximal diagonal domain.

Unique factorization makes \(U\) a bijection of orthonormal bases.
On an occupation vector, second quantization has eigenvalue
\(\prod_p p^{-sr_p}\), which proves the operator equality.
Its trace-class domain follows from Proposition~\ref{prop:prime-product}.
For a finite set of primes \(F\), summing all occupations gives
\(\prod_{p\in F}(1-p^{-s})^{-1}\).  The projections onto configurations
supported in \(F\) tend to the identity.  For \(\operatorname{Re}s>1\),
the omitted absolute diagonal sum tends to zero, so the convergence is
in trace norm.  On \(\mathscr K\), the finite compressions converge in
\(\mathcal S_1\).  Their determinants converge to
\(\prod_p(1-p^{-s})\), the Fredholm determinant of the diagonal operator.
Both products converge locally uniformly.  This proves
the trace--determinant identity, including its analytic meaning and normalization.
\end{proof}

\subsection{Regularization and its exact boundary}

When \(0<\operatorname{Re}s\leq1\), the prime operator remains compact.
The divergent first terms of its logarithmic determinant can be removed.
For \(m\geq1\), define the entire function
\[
 E_m(z)=(1-z)\exp\left(\sum_{k=1}^{m-1}\frac{z^k}{k}\right),
\]
where an empty sum is zero.  Since \(E_m(z)-1=O(z^m)\) at zero,
\(E_m(K)-I\) is trace class for the diagonal \(K\in\mathcal S_m\)
used here.  Define
\[
 \det_m(I-K)=\det\bigl(E_m(K)\bigr).
\]
This fixes the regularization, including every exponential factor.

\begin{theorem}[Regularized prime determinants]
\label{thm:prime-regularization}
For every integer \(m\geq1\),
\[
 \Delta_m(s):=\det_m(I-K(s))
 =\prod_p E_m(p^{-s})
\]
is nonzero and holomorphic on \(\operatorname{Re}s>1/m\).  It has the
unique holomorphic logarithm tending to zero along the positive real axis,
\[
 L_m(s):=\log \Delta_m(s)
 =-\sum_p\sum_{k\geq m}\frac{p^{-ks}}{k}.
\]
If \(P(z)=\sum_p p^{-z}\) for \(\operatorname{Re}z>1\), then
\[
 \Delta_m(s)=\zeta(s)^{-1}
 \exp\left(\sum_{k=1}^{m-1}\frac{P(ks)}{k}\right),
 \qquad \operatorname{Re}s>1.
\]
For \(\sigma=\operatorname{Re}s>1/m\) and integer \(R\geq2\),
\[
 \left|L_m(s)-\sum_{p\leq R}\log E_m(p^{-s})\right|
 \leq\frac{R^{1-m\sigma}}
 {m(m\sigma-1)(1-2^{-\sigma})}.
\]
On the same initial half-plane, the full logarithmic derivative is
\[
 L_m'(s)=\sum_p(\log p)\frac{p^{-ms}}{1-p^{-s}}.
\]
The logarithms in the tail estimate use the power series
\(\log E_m(z)=-\sum_{k\geq m}z^k/k\) for \(|z|<1\).
\end{theorem}

\begin{proof}
For \(|z|\leq2^{-\sigma}<1\), the defining series gives
\[
 |\log E_m(z)|\leq\frac{|z|^m}{m(1-2^{-\sigma})}.
\]
Summation over primes proves absolute and locally uniform convergence
of the logarithmic series.  Exponentiation gives a nonzero holomorphic
product.  The trace-class diagonal \(E_m(K(s))-I\) has the same Fredholm
product, so this is the stated operator determinant.  Its limit at
\(+\infty\) is \(1\), and the logarithm has limit zero.
On \(\operatorname{Re}s>1\), add back the omitted terms and apply
Theorem~\ref{thm:euler-fredholm}.  Finally,
\[
 \sum_{p>R}p^{-m\sigma}
 \leq\sum_{n>R}n^{-m\sigma}
 \leq\int_R^\infty x^{-m\sigma}\,dx
\]
proves the stated tail bound.  On a smaller closed half-plane within
this domain, differentiation is justified by
\(\sum_{n\geq2}(\log n)n^{-m\sigma}<\infty\).
Differentiating each summand and summing its geometric series gives the
stated full logarithmic derivative.
\end{proof}

The two-prime example already distinguishes the normalizations:
\[
 \det_2(I-K_{\{2,3\}}(s))
 =\det(I-K_{\{2,3\}}(s))\exp(2^{-s}+3^{-s}).
\]
At \(s=1\), its value is \(e^{5/6}/3\).  A constant cannot remove this
parameter-dependent correction.

\begin{theorem}[Fractional zero at the convergence boundary]
\label{thm:fractional-zero}
For \(m\geq2\), there is a disk about \(a_m=1/m\) and a nonvanishing
holomorphic function \(G_m\) on that disk such that
\[
 \Delta_m(s)^m=(ms-1)G_m(s),\qquad \operatorname{Re}s>a_m.
\]
Moreover, \(G_m(a_m)>0\).  Thus \(\Delta_m\) has no single-valued meromorphic
continuation to a neighborhood of \(a_m\).  On the real axis,
\[
 \Delta_m(s)\sim C_m(ms-1)^{1/m}\quad(s\downarrow a_m),\qquad C_m>0.
\]
\end{theorem}

\begin{proof}
Put
\[
 R_0(z)=\sum_{k\geq2}\frac{P(kz)}k,\qquad
 B_m(s)=\sum_{k\geq m+1}\frac{P(ks)}k.
\]
The estimate in Theorem~\ref{thm:prime-regularization} shows that these
functions are holomorphic on \(\operatorname{Re}z>1/2\) and
\(\operatorname{Re}s>1/(m+1)\), respectively.
For \(\operatorname{Re}z>1\),
\[
 \log\zeta(z)=P(z)+R_0(z).
\]
The residue of \(\zeta\) at \(1\) is one.  One can see its normalization
directly from
\[
 \zeta(z)-\frac1{z-1}
 =\sum_{n\geq1}\left(n^{-z}-\int_n^{n+1}x^{-z}\,dx\right).
\]
The right side converges locally uniformly for \(\operatorname{Re}z>0\),
because its summands are \(O(n^{-\operatorname{Re}z-1})\) on compact
sets.  Consequently \(Q(z)=(z-1)\zeta(z)\) is holomorphic near \(1\),
with \(Q(1)=1\), and is nonzero on a sufficiently small disk.

The logarithmic series of Theorem~\ref{thm:prime-regularization} gives, for
\(\operatorname{Re}s>a_m\),
\[
 \Delta_m(s)^m
 =\zeta(ms)^{-1}\exp\bigl(R_0(ms)-mB_m(s)\bigr)
 =(ms-1)\frac{\exp\bigl(R_0(ms)-mB_m(s)\bigr)}{Q(ms)}.
\]
The last quotient is the required \(G_m\).  It is holomorphic and nonzero
near \(a_m\), with
\[
 G_m(a_m)=\exp\bigl(R_0(1)-mB_m(a_m)\bigr)>0.
\]
If \(\Delta_m\) were meromorphic there, the order of its \(m\)-th power
would be divisible by \(m\).  The displayed simple zero contradicts this.
Since every factor is positive for real \(s>a_m\), the asymptotic follows
with \(C_m=\exp(R_0(1)/m-B_m(a_m))\).
\end{proof}

The fractional zero suggests the coordinate \(t^m=ms-1\), in which
the determinant has an ordinary zero.  The next section constructs its
continuation surface and the comparison maps for different regularization
orders from the convergent remainder functions.

\section{Branched continuation of the prime determinant}
\label{sec:prime-continuation}

A fractional zero can become an ordinary zero after changing the
coordinate on the domain.  The square root of \(s-a\) returns with the
opposite sign after one counterclockwise circuit around \(a\).
On the coordinate surface \(s=a+t^2\), it is the single-valued function
\(t\).  For the prime determinant, the analogous construction will require
infinitely many different ramification orders.

Fix an integer \(m\geq2\) and retain precisely \(E_m\), \(L_m\), and
\(\Delta_m\) from Theorem~\ref{thm:prime-regularization}.
Thus the initial domain is \(\operatorname{Re}s>1/m\),
\(p^{-s}=\exp(-s\log p)\) uses the real logarithm of \(p\), and
\(L_m(s)\to0\) as real \(s\to+\infty\).
Our base point is \(s_*=2\), with the positive value \(\Delta_m(2)\).
The continuation concerns this fixed holomorphic germ; the summability
domain of the operator in Theorem~\ref{thm:euler-fredholm} stays fixed.

We work over \(\mathbb C\) and use its usual topology.  A meromorphic
function is locally a quotient of holomorphic functions.  At a point \(a\),
its order is the integer \(v\) in \(f(s)=(s-a)^v u(s)\), where \(u\)
is holomorphic and nonzero.  Positive orders are zeros and negative
orders are poles.  The divisor records these points with their integer
orders; its support is the set of points with nonzero order.
A covering map is locally a disjoint union of copies
of an open set.  A branched covering also permits local coordinates
\(s=a+t^q\), with \(q\) a positive integer.  We will construct these
coordinates and their gluing maps explicitly.

We use convergent complex power series and integration along piecewise
smooth paths.  The local logarithms and residue formula needed for the
continuation maps are developed next.

\subsection{Logarithms, residues, and zeta factors}

\subsubsection{Local logarithms and integrals}

A locally uniform limit of holomorphic functions is holomorphic.
For completeness, on a circle contained in the domain, uniform convergence
allows passage to the limit in the Cauchy integral formula.
Expanding \((z-s)^{-1}\) as a geometric series on a smaller concentric
disk gives a convergent power series for the limit.
The same argument gives locally uniform convergence of derivatives there.
The integral formula itself follows by integrating the power-series
primitive on small disks and subdividing the region between the circle
and a smaller circle about \(s\).  Internal edges cancel, while the
smaller circle contributes \(2\pi i f(s)\).

If \(u(a)\ne0\), shrink a disk about \(a\) until
\(|u(s)/u(a)-1|<1\).  The power series for \(\log(1+z)\), together
with one logarithm of \(u(a)\), gives a holomorphic logarithm of \(u\).
Two logarithms on a connected overlap differ by \(2\pi i n\),
\(n\in\mathbb Z\), because their exponentials agree and the difference
is continuous with discrete values.  Roots of a nonvanishing function
are obtained by dividing such a logarithm by the required integer.

For a closed path \(\gamma\) avoiding \(a\), define
\[
 \operatorname{wind}(\gamma,a)
 =\frac1{2\pi i}\int_\gamma\frac{ds}{s-a}.
\]
To see that this is an integer, continue a local logarithm of \(s-a\)
along finitely many small disks covering the path.  At the endpoint its
change is \(2\pi i n\).  Integrating its derivative gives the formula.
Subdividing a homotopy into small rectangles shows that the integral
of a holomorphic one-form is unchanged under path homotopy with fixed
endpoints: each rectangle lies in a disk with a power-series primitive.

We need a related residue formula.  Suppose a one-form on a half-plane
is holomorphic except at a locally finite set of simple poles \(a\),
with residues \(r_a\).  A closed path lies inside a compact convex polygon
in that half-plane whose boundary avoids the poles.
There are only finitely many poles in the polygon.
Subtract \(\sum_a r_a\,ds/(s-a)\).  The remaining form is holomorphic
on a slightly smaller polygon containing the path and has zero integral
around closed paths by subdivision.  It follows that
\[
 \int_\gamma\omega
 =2\pi i\sum_a r_a\operatorname{wind}(\gamma,a).
\]
Only finitely many terms contribute.  These elementary integral identities
will determine the entire continuation space.

\subsubsection{Continuation of the ordinary zeta function}

Write \(\Hp_c=\{s\in\mathbb C:\Re s>c\}\), and \(\Hp=\Hp_0\).
For \(\Re z>1\), let \(\zeta(z)=\sum_{n\geq1}n^{-z}\).
Finite geometric products and unique factorization give
\[
 \zeta(z)=\prod_p(1-p^{-z})^{-1},\qquad
 Z(z):=\log\zeta(z)=\sum_p\sum_{j\geq1}\frac{p^{-jz}}j.
\]
Both limits are locally uniform and absolute on \(\Hp_1\).
For example, all omitted integers in a finite prime product are bounded
in absolute sum by the tail of \(\sum n^{-\Re z}\).
The absolute logarithmic sum is bounded by
\((1-2^{-\Re z})^{-1}\sum_{n\geq2}n^{-\Re z}\).
Its exponential is the product, so \(\zeta\) is nonzero on \(\Hp_1\).
The branch \(Z\) tends to zero on the positive real axis.

The subtraction formula in the proof of
Theorem~\ref{thm:fractional-zero} continues \(\zeta\) meromorphically
to \(\Hp\), with its only pole at \(1\), simple and of residue one.
Its locally uniform derivative bound proves holomorphy of the subtracted
series throughout \(\Hp\).  This construction is sufficient for every
zeta factor used below.

Zeros are isolated.  A holomorphic function that is not identically zero
has a first nonzero Taylor coefficient at each zero, which factors it as
\((z-a)^v u(z)\) with \(u(a)\ne0\).  An accumulation of zeros would
make every coefficient vanish and then, by overlapping disks, make the
function identically zero.  The zeta function is nonzero on \(\Hp_1\),
so this cannot occur.

We will also need the absence of zeros on the positive real axis.
For \(x>0\), the paired series
\[
 \eta(x)=\sum_{r\geq1}\bigl((2r-1)^{-x}-(2r)^{-x}\bigr)
\]
converges and is strictly positive.  Its complex analogue converges
locally uniformly on \(\Hp\) by the same derivative estimate.
Splitting even and odd terms for \(\Re z>1\), and then continuing,
gives \(\eta(z)=(1-2^{1-z})\zeta(z)\).
Thus \(\zeta(x)<0\) for \(0<x<1\), while \(\zeta(x)>0\) for \(x>1\).

\subsection{The first continuation surface}

The proof of Theorem~\ref{thm:fractional-zero} isolates the first
divergent part without continuing the prime series itself.
Retain its functions \(P\), \(R_0\), and \(B_m\), holomorphic on
\(\Hp_1\), \(\Hp_{1/2}\), and \(\Hp_{1/(m+1)}\), respectively.
The same calculation gives, on \(\Hp_{1/m}\),
\[
 mL_m(s)=-P(ms)-mB_m(s)
 =-Z(ms)+R_0(ms)-mB_m(s).
\]
Consequently the single-valued meromorphic function
\[
 F_m(s)=\frac{\exp(R_0(ms)-mB_m(s))}{\zeta(ms)}
 \quad\text{on }U_m:=\Hp_{1/(m+1)}
\]
extends \(\Delta_m^m\).  The stated domain follows because
\(m/(m+1)>1/2\).  The numerator is holomorphic and nonzero throughout it.
The function \(F_m\) has a simple zero at \(a_m=1/m\).
At a point \(\rho/m\) where \(\zeta(\rho)=0\) has multiplicity \(v\)
and \(\Re\rho>m/(m+1)\), it has a pole of order \(v\).
There are no other zeros or poles on \(U_m\).

\begin{theorem}[Minimal first surface]
\label{thm:first-surface}
Let \(U_m^\circ\) be \(U_m\) with the divisor support of \(F_m\)
removed.  The equation
\[
 Y_m=\{(s,y):s\in U_m^\circ,\ y\in\mathbb C^*,\ y^m=F_m(s)\}
\]
defines a connected \(m\)-sheeted holomorphic cover.
It carries the continuation of \(\Delta_m\) as the function \(y\),
with the initial sheet specified by \(y(2)=\Delta_m(2)>0\).
It admits a meromorphic branched filling over every divisor point.
If \(v=\ord_a F_m\ne0\) and \(d=\gcd(m,|v|)\), the filling has
\(d\) points over \(a\).  Each has ramification index \(m/d\), and
\(y\) has order \(v/d\) there.
At \(a_m\), this is the minimal local \(m\)-fold cover, and \(y\)
has a simple zero.
\end{theorem}

\begin{proof}
On a small disk avoiding the divisor, choose \(\log F_m\).
The \(m\) functions
\(y_j=\exp((\log F_m+2\pi i j)/m)\), \(0\leq j<m\), are disjoint
local sheets.  A circuit about the simple zero at \(a_m\) permutes them
by \(j\mapsto j+1\).  Thus all sheets are connected by continuation.

Near a divisor point, write \(F_m(s)=(s-a)^v u(s)\) and choose
\(h^m=u\) with \(h\) holomorphic and nonzero.  The charts are
\[
 s=a+t^{m/d},\qquad
 y=\xi\,t^{v/d}h(a+t^{m/d}),\qquad \xi^m=1.
\]
Negative powers of \(t\) mean a meromorphic pole.  Changing \(t\) by
an \((m/d)\)-th root of unity identifies exactly \(m/d\) choices of
\(\xi\), because \(v/d\) and \(m/d\) are coprime.  Hence there are
\(d\) charts.  They parametrize all the punctured sheets and fill their
ends with smooth disk coordinates.  On overlaps their changes of
coordinate are the roots of unity just described, so they glue.

At \(a_m\), write \(F_m(s)=(ms-1)G_m(s)\), where
\[
 G_m(s)=\frac{\exp(R_0(ms)-mB_m(s))}{(ms-1)\zeta(ms)}.
\]
The denominator is nonzero near \(a_m\) and equals one there.
Thus \(G_m(a_m)>0\).  The coordinate \(t^m=ms-1\) gives
\(y=t\,G_m((1+t^m)/m)^{1/m}\), using the root positive at \(t=0\).
One circuit multiplies \(y\) by \(e^{2\pi i/m}\).
On any local cover \(s-a_m=t^e\) where \(y\) is meromorphic,
its order must be \(e/m\), an integer.  Therefore \(m\) divides \(e\).
Our chart attains this bound.
\end{proof}

For \(m=2\), the two sheets meet smoothly above \(1/2\), and one
circuit changes the sign.  At a zero of \(\zeta(2s)\) of order two,
the same construction gives two unramified simple poles.
A simple zero of \(\zeta(2s)\) would instead give one ramified simple
pole upstairs.  These are conditional descriptions of any such zeros
in \(U_2\), not assertions that they exist there.

\subsection{Finite extraction and continuation through the positive half-plane}

The preceding surface uses the full domain of the two displayed remainder
series.  To go farther left, remove finitely many more logarithmic factors.
The coefficients determine every branch.

In the extraction and the cover construction that follows, allow any
integer \(m\geq1\); order one will recover the reciprocal zeta function.
Define the integer \(\mu(n)\) as zero when a prime square divides \(n\),
and as \((-1)^r\) when \(n\) is a product of \(r\) distinct primes.
Set \(\mu(1)=1\).  Choosing subsets of the prime divisors proves
\(\sum_{d\mid n}\mu(d)=0\) for \(n>1\), and the sum is one for \(n=1\).
For \(n\geq1\), put
\[
 c_m(n)=\frac1n\sum_{\substack{k\mid n\\ k\geq m}}\mu(n/k).
\]
Then \(c_m(n)=0\) for \(n<m\), \(c_m(m)=1/m\), and
\(|c_m(n)|\leq1\).  Summing over divisors in the opposite order gives
\[
 \sum_{d\mid n}d\,c_m(d)=\mathbf1_{n\geq m}.
\]
Indeed, the contribution for a fixed \(k\mid n\), \(k\geq m\), is
\(\sum_{e\mid n/k}\mu(e)\), so only \(k=n\) survives.

\begin{theorem}[Finite logarithmic extraction]
\label{thm:extraction}\label{mono:prop:extraction}
On \(\Hp_{1/m}\),
\[
 L_m(s)=-\sum_{n\geq m}c_m(n)Z(ns).
\]
For \(N\geq m\), the remainder
\[
 T_{m,N}(s)=-\sum_{n>N}c_m(n)Z(ns)
\]
is single-valued and holomorphic on \(U_N=\Hp_{1/(N+1)}\).
On every open disk in \(U_N\) avoiding zeros and poles of
\(\zeta(ns)\), \(m\leq n\leq N\), the formula
\[
 f_{m,N}(s)=\exp\left(T_{m,N}(s)
       -\sum_{n=m}^N c_m(n)\ell_n(s)\right),
 \qquad e^{\ell_n(s)}=\zeta(ns),
\]
gives a local branch of the continued determinant, when its logarithms
are continued from the initial half-plane.
The formulas for successive \(N\) agree with these continued choices.
\end{theorem}

\begin{proof}
For \(x>1\) and \(\Re z=x\),
\[
 |Z(z)|\leq\frac1{1-2^{-x}}\sum_{r\geq2}r^{-x}
 \leq\frac{2^{-x}+2^{1-x}/(x-1)}{1-2^{-x}}.
\]
For \(x\geq2\), this is at most \(4\,2^{-x}\).
On a compact subset of \(U_N\), all \(ns\), \(n>N\), lie in
\(\Hp_1\).  The large \(n\) terms therefore have a geometric bound,
which proves local uniform convergence of \(T_{m,N}\) and its derivatives.

On \(\Hp_{1/m}\), expand each \(Z(ns)\) into its prime series.
Absolute convergence follows from the same bound, so sums can be reordered.
The coefficient of \(p^{-js}\) becomes
\[
 -\sum_{n\mid j}\frac{c_m(n)}{j/n}
 =-\frac{\mathbf1_{j\geq m}}j.
\]
This is the defining coefficient of \(L_m\).
The local formula follows by keeping its first \(N\) terms separate.
On \(U_N\),
\[
 T_{m,N}=T_{m,N+1}-c_m(N+1)Z((N+1)s).
\]
The last logarithm has its original single-valued branch there.
For \(m=1\), divisor cancellation gives
\(c_1(n)=\mathbf1_{n=1}\), so the same identity reduces to
\(L_1=-Z\).  The compatibility therefore holds for every positive
integer order.
\end{proof}

For example, \(c_2(n)=-\mu(n)/n\) for \(n>1\).
The first coefficients are
\[
 c_2(2)=\frac12,\quad c_2(3)=\frac13,\quad c_2(4)=0,
 \quad c_2(5)=\frac15,\quad c_2(6)=-\frac16.
\]
The first continuation takes a square root of an inverse zeta factor.
The next takes a cube root as well.  The negative coefficient at six
reverses the contribution of the pole of \(\zeta(6s)\).

\subsubsection{The logarithmic derivative and coincident divisors}

Let \(S_m\subset\Hp\) consist of points where \(\zeta(ns)\) has a
zero or pole for some \(n\geq m\) with \(c_m(n)\ne0\).
This set is locally finite.  On a compact set with \(\Re s\geq\epsilon>0\),
only \(n\leq1/\epsilon\) can contribute, since \(\zeta\) is nonzero
and holomorphic on \(\Hp_1\).  Each of these finitely many factors
has finitely many zeros and poles on the compact set.
Put \(\Omega_m=\Hp\setminus S_m\).

The locally finite punctures do not disconnect the half-plane.
A polygonal path between two points meets only finitely many punctures
in a compact neighborhood, and small detours avoid them.
Hence \(\Omega_m\) is path connected.

On \(\Omega_m\), define the holomorphic one-form
\[
 \omega_m=dT_{m,N}
 -\sum_{n=m}^N c_m(n)\frac{n\zeta'(ns)}{\zeta(ns)}\dd
 \quad\text{where }s\in U_N.
\]
Theorem~\ref{thm:extraction} proves independence of \(N\).
For every \(a\in\Hp\), define the rational number
\[
 \alpha_m(a)=-\sum_{n\geq m}c_m(n)\ord_a\zeta(ns).
\]
Only finitely many terms occur.  This is the residue of the meromorphic
extension of \(\omega_m\), with value zero wherever that extension is
holomorphic.  The signs include a zeta pole with order
\(-1\).  Coincident zeros and poles contribute to the same sum.
Each logarithmic derivative has a simple pole, so if this sum vanishes,
\(\omega_m\) extends holomorphically at \(a\).
For a loop \(\gamma\) in \(\Omega_m\), set
\[
 \chi_m(\gamma)=\exp\left(\int_\gamma\omega_m\right)
 =\exp\left(2\pi i\sum_{a\in S_m}
       \alpha_m(a)\operatorname{wind}(\gamma,a)\right).
\]
This is a homomorphism from based loops modulo homotopy to
the unit complex numbers.  Every value is a root of unity.
It is the monodromy character of the determinant.

\subsection{The minimal cover and the determinant line}

A Riemann surface is a Hausdorff space with a countable atlas of open
subsets of \(\mathbb C\) and holomorphic coordinate changes.
A pointed germ here means a holomorphic function near \(2\), with
two such functions identified if they agree on a smaller neighborhood.
A torsor for a group is a set on which the group acts so that exactly
one group element carries any specified point to any other.
For loops based at \(2\), multiplication means concatenation, and
reversing a path gives the inverse.  The kernel of \(\chi_m\) consists
of loops with character one.

\begin{theorem}[Continuation on a minimal cover]
\label{thm:minimal-cover}\label{mono:prop:local}
There is a connected holomorphic covering
\(\pi:\mathcal E_m\to\Omega_m\) and a nonzero holomorphic function
\(\mathcal D_m\) on it, with the following properties.
The pointed germ at \(2\) is \(\Delta_m\), and
\[
 \frac{d\mathcal D_m}{\mathcal D_m}=\pi^*\omega_m.
\]
The fiber of \(\pi\) is a torsor for \(\operatorname{im}\chi_m\).
Every pointed connected cover on which that germ becomes single-valued
maps uniquely to \(\mathcal E_m\) over \(\Omega_m\), preserving the germ.
Thus \(\mathcal E_m\) has the fewest sheets among such covers.
It has a meromorphic branched filling over every point of \(S_m\).
If \(\alpha_m(a)=b/q\) in lowest terms with \(q>0\), each local end
has ramification index \(q\), and \(\mathcal D_m\) has order \(b\).
For \(\alpha_m(a)=0\), take \(b=0,q=1\).
Every meromorphic realization with local ramification index \(e\)
requires \(q\mid e\).
\end{theorem}

\begin{proof}
Consider paths \(\gamma\) in \(\Omega_m\) from \(2\).
Two paths are equivalent when they have the same endpoint and
\(\exp(\int_\gamma\omega_m-\int_\eta\omega_m)=1\).
Write \([\gamma]\) for the class, and set
\[
 \pi[\gamma]=\gamma(1),\qquad
 \mathcal D_m[\gamma]=\Delta_m(2)
       \exp\left(\int_\gamma\omega_m\right).
\]
The equivalence relation makes the function well-defined.
For a small disk \(V\subset\Omega_m\) about \(\gamma(1)\), append
paths inside \(V\).  A primitive of \(\omega_m\) on \(V\) shows
that the resulting classes form a copy of \(V\).
Two such copies are either disjoint or identical over a connected overlap,
with change of coordinates equal to the identity on the base.
These charts define the covering and its complex structure.
They separate distinct points in a fiber, so the space is Hausdorff.
There are at most countably many classes in a fiber, because finite sums
of rational residues give a countable monodromy image.
A countable disk basis on the base therefore gives a countable atlas.
Every class is joined to the base class by the initial portions of its
representing path.  Thus the cover is connected.

The displayed exponential is holomorphic in each chart, with the stated
logarithmic derivative.  On the initial half-plane it agrees with
\(\Delta_m\), since the logarithmic derivatives and the values at \(2\)
agree.  Appending loops gives the simply transitive action of
\(\operatorname{im}\chi_m\) on each fiber.

Let \(X\to\Omega_m\) be another pointed connected cover carrying the
continued germ.  A path in \(X\) from its base point projects to a path
\(\gamma\).  Its endpoint value is forced by the exponential formula.
Two paths ending at the same point therefore have trivial exponential
period.  Sending that point to \([\gamma]\) is well-defined, holomorphic,
and unique.  It is onto because every path on the base lifts to \(X\).
This proves minimality.

Near a puncture, the Laurent calculation gives
\[
 \omega_m=\frac bq\frac{ds}{s-a}+dh(s)
\]
for a holomorphic \(h\) on a disk.  Each branch is consequently
\(C(s-a)^{b/q}e^{h(s)}\), with \(C\ne0\).
The punctured sheets form cycles of length \(q\), since \(b,q\) are
coprime.  Fill each cycle with the coordinate \(s=a+t^q\).
The function becomes \(Ct^b e^{h(a+t^q)}\), which is meromorphic.
Different choices of the coordinate differ by a \(q\)-th root of unity.
This gives the compatible local filling and the stated order.
On a chart \(s-a=t^e\), a meromorphic realization would have integer
order \(eb/q\); coprimality of \(b,q\) forces \(q\mid e\).
\end{proof}

The formula resolves cancellations before assigning ramification.
For example, local exponents \(1/2\) and \(-1/2\) give a nonzero
holomorphic branch, not a double cover.  A total exponent \(2/3\)
gives a triple cover with a double zero.  An exponent \(-2/3\) gives
the same cover with a double pole.  The denominator of the reduced total
exponent, rather than any individual factor, is the minimal local degree.

\subsubsection{A finite description on each half-plane}

The construction has finite equations on every \(U_N\).
Let \(M_N\) be the least common multiple of the denominators of
\(c_m(m),\ldots,c_m(N)\), in lowest terms.  Define
\[
 A_{m,N}(s)=\exp(M_N T_{m,N}(s))
       \prod_{n=m}^N\zeta(ns)^{-M_Nc_m(n)}.
\]
This is meromorphic on \(U_N\), since all powers are integers.
Over \(U_N\setminus S_m\), the equation
\(y^{M_N}=A_{m,N}(s)\) is connected and has exactly \(M_N\) sheets.
To prove this, first note the identity
\[
 \alpha_m(1/n)=c_m(n),\qquad n\geq m.
\]
Only \(\zeta(ns)\) has a pole at \(1/n\); all other scaled zeta
values there are nonzero finite values on the positive real axis.
For \(m\leq n\leq N\), these points lie in \(U_N\).
Small loops about them generate the subgroup of roots of unity generated
by \(\exp(2\pi i c_m(n))\).  Its order is \(M_N\): a positive integer
annihilates every exponent modulo \(\mathbb Z\) exactly when it is
divisible by every reduced denominator.  Thus continuation reaches all
the roots of the polynomial equation.  Conversely every residue on
\(U_N\) has denominator dividing \(M_N\), so no further sheets occur.
The path construction identifies this equation, with base value
\(y(2)=\Delta_m(2)\), with the minimal continuation using paths in
that half-plane.

For order two, the first degrees are
\[
 \begin{array}{c|ccccc}
 N&2&3&4&5&6\\ \hline
 M_N&2&6&6&30&30
 \end{array}
\]
The point \(1/4\) introduces no sheet, since \(c_2(4)=0\).
At \(1/6\), the exponent \(-1/6\) gives a simple pole after the
sixfold local pullback.  Both facts occur within the degree-thirty surface
over \(U_6\).  For order three, the exponents \(1/3\) and \(1/4\)
already force degree twelve over \(U_4\).

Since \(M_N\mid M_{N+1}\), the remainder identity implies
\[
 A_{m,N+1}|_{U_N}=A_{m,N}^{\,M_{N+1}/M_N}.
\]
The transition between the pointed components is exactly
\((s,y)\mapsto(s,y)\).  Its image is the pointed component over
\(U_N\) inside the next surface.  The maps also extend over filled
divisor points, since both use the same reduced residue and local coordinate.

Every path in \(\Omega_m\) is compact and has positive minimum real
part.  It therefore lies in some \(U_N\).  The union of these compatible
pointed surfaces is \(\mathcal E_m\).  This proves continuation through
the whole positive half-plane with its described branched fillings.
It is the full domain supplied by this exhaustion of convergent
remainders.  No assertion of maximal continuation outside \(\Hp\) is
needed or proved.

\subsubsection{Branches as a section of a line bundle}

There is also a description without enlarging the base.
Choose small disks \(V_i\) covering \(\Omega_m\), and continue the
initial determinant to a branch \(f_i\) on each disk along a chosen path.
On a connected component of \(V_i\cap V_j\), put
\[
 g_{ij}=f_i/f_j.
\]
Both branches have logarithmic derivative \(\omega_m\), so \(g_{ij}\)
is constant.  It is a value of \(\chi_m\), and
\(g_{ij}g_{jk}=g_{ik}\) on triple overlaps.
Identify \((s,v_j)\in V_j\times\mathbb C\) with
\((s,v_i)\in V_i\times\mathbb C\) when \(v_i=g_{ij}v_j\).
The cocycle identity makes the quotient a holomorphic line bundle
\(\mathcal L_m\).  The values \(f_i\) define a nonzero holomorphic
section of it.  All bundle and section maps have now been specified.

Because the transition constants have zero derivative, ordinary
differentiation of local coordinates defines a flat connection \(\nabla\)
on \(\mathcal L_m\).  Flat means that continuing a locally constant
coordinate around a sufficiently small loop returns the same coordinate.
The determinant section satisfies \(\nabla f=\omega_m f\).
Its nonvanishing trivializes the underlying holomorphic line bundle.
In that trivialization, the connection is \(d+\omega_m\).
Its parallel transport around \(\gamma\) multiplies by
\(\chi_m(\gamma)^{-1}\).  Thus the continuation information lies in the
specified connection as well as the section; it is not a claim that the
underlying holomorphic line bundle is nontrivial.
The branched charts of Theorem~\ref{thm:minimal-cover} describe its
meromorphic extension after the required local pullback.

\subsection{Changing the regularization order}

The order of regularization changes the function.  The comparison must
therefore include a multiplier and the space on which it is defined.
For \(n>m\), define on \(\Hp_{1/m}\)
\[
 H_{m,n}(s)=\exp\left(\sum_{k=m}^{n-1}\frac{P(ks)}k\right).
\]
Subtracting the absolutely convergent defining series gives
\[
 \Delta_n=H_{m,n}\Delta_m,\qquad
 H_{m,n}H_{n,r}=H_{m,r}\quad(m<n<r).
\]
For instance, \(H_{2,3}=\exp(P(2s)/2)\).
At \(s=1/2\), its continued local exponent is \(-1/2\), canceling
the exponent \(1/2\) of \(\Delta_2\).  The resulting \(\Delta_3\)
is already holomorphic and nonzero there by its initial definition.

\begin{theorem}[Comparison on common covers]
\label{thm:order-comparison}
On \(\Omega=\Omega_m\cap\Omega_n\), the multiplier \(H_{m,n}\)
continues with logarithmic derivative \(\omega_n-\omega_m\) and
character \(\chi_n\chi_m^{-1}\).  On the pointed cover defined by
\(\ker\chi_m\cap\ker\chi_n\), all three functions are single-valued
and \(\mathcal D_n=H_{m,n}\mathcal D_m\).
At a divisor point, the multiplier has local exponent
\(\alpha_n-\alpha_m\).  The comparisons preserve their cocycle identity
on common covers and in the corresponding line bundles.

Multiplication descends to a map from the minimal \(m\)-cover to the
minimal \(n\)-cover precisely when
\(\ker\chi_m\subseteq\ker\chi_n\), with characters restricted to
\(\Omega\).  Here each minimal cover uses paths in \(\Omega\) from
the fixed base point.  Such descent is not automatic.
\end{theorem}

\begin{proof}
The initial identity gives the difference of logarithmic derivatives and
fixes the multiplier at \(2\).  Path integration therefore defines its
unique continuation.  The exponential period is the quotient of the two
characters.  It is trivial on their common kernel, so multiplication is
well-defined on the stated common cover.  The local exponent follows by
subtracting residues.  For three orders, the logarithmic derivatives
telescope and the base values multiply correctly.  Uniqueness of the
continued germ proves the cocycle identity.

For the line bundles, choose the same continuation paths on a common disk
cover.  The multiplier has transition constants
\(g_{ij}^{(n)}/g_{ij}^{(m)}\), so it is a nonzero section of
\(\mathcal L_n\otimes\mathcal L_m^{-1}\).
Multiplying local coordinates gives an isomorphism of holomorphic line
bundles carrying the determinant section to the other one.  The difference
of the two flat connections is the specified logarithmic derivative.

A point of the minimal \(m\)-cover identifies paths whose quotient loop
lies in \(\ker\chi_m\).  Sending it to the corresponding path class
in the \(n\)-cover is well-defined exactly when every such loop lies
in \(\ker\chi_n\).  This proves the descent criterion in both directions.
\end{proof}

For a concrete failure, compare orders \(2\) and \(3\) around \(s=1/4\).
The coefficient formulas give \(c_2(4)=0\) and \(c_3(4)=1/4\).
All other scaled zeta factors are finite and nonzero at that positive
real point.  Thus \(\chi_2=1\) and \(\chi_3=i\) on its small positive
loop.  The order-three branch cannot be a function on the minimal
order-two cover near this loop.  The common cover in the theorem supplies
the required comparison without discarding this monodromy.

\subsection{Why the number of sheets is unbounded}

The finite equations determine the whole sheet group, including its
prime-power orders.  For a prime \(p\), let \(e_p\) be the least
positive integer such that \(p^{e_p}\geq m\).  Equivalently,
\(p^{e_p-1}<m\leq p^{e_p}\).
For a positive integer \(d\), let \(\mu_d\subset\mathbb C^*\) be the
group of \(d\)-th roots of unity.

\begin{theorem}[Exact monodromy group]
\label{thm:monodromy-group}
For \(m\geq2\), the image of the continuation character is
\[
 \operatorname{im}\chi_m
 =\bigcup_{N\geq m}\mu_{M_N}
 =\{z\in\mathbb C^*:z\text{ has finite order }d,
       \ v_p(d)\leq e_p\text{ for every prime }p\}.
\]
Here \(v_p(d)\) is the exponent of \(p\) in the prime factorization of
\(d\).  As an abstract group, this is
\(\bigoplus_p\mathbb Z/p^{e_p}\mathbb Z\): only finitely many
coordinates are nonzero in each element.
\end{theorem}

\begin{proof}
Every loop lies in some \(U_N\), whose monodromy image is exactly
\(\mu_{M_N}\).  This proves the first equality.
For \(n>1\), divisor cancellation rewrites the coefficient as
\[
 c_m(n)=-\frac1n\sum_{\substack{k\mid n\\k<m}}\mu(n/k).
\]
If \(v_p(n)>e_p\), then \(k<m\) implies \(v_p(k)\leq e_p-1\).
Thus \(p^2\mid n/k\) in every term, so \(c_m(n)=0\).
Every nonzero coefficient consequently has denominator with
\(p\)-valuation at most \(e_p\).

Conversely, take \(n=p^{e_p}\).  Its only divisor at least \(m\) is
\(n\) itself, so \(c_m(n)=1/p^{e_p}\).
The loop about \(1/p^{e_p}\) gives a generator of \(\mu_{p^{e_p}}\).
These loops for finitely many distinct primes generate all roots with
order dividing their product.  Indeed, cyclic groups of coprime orders
intersect trivially and their product has the product order.
This proves the second equality and the direct-sum description.
\end{proof}

For \(m=2\), every \(e_p=1\), so the possible orders are precisely
the squarefree positive integers.  For \(m=3\), one has \(e_2=2\),
while \(e_p=1\) for every odd prime.  The allowed orders divide four
times an odd squarefree integer.  Increasing the regularization can change
this group, as well as the local determinant values.

\begin{theorem}[Prime ramification and the boundary at zero]
\label{thm:infinite-degree}
Fix \(m\geq2\).  For every prime \(q\geq m\),
\[
 \alpha_m(1/q)=1/q.
\]
Hence the minimal continuation cover over \(\Omega_m\) has infinite
degree.  No finite-degree branched cover of the whole positive half-plane
carries the continued determinant as a meromorphic function.
Nor can one finite-degree branched cover of a neighborhood of zero carry
all these continued germs at the positive points \(1/q\).
\end{theorem}

\begin{proof}
At \(s=1/q\), the factor \(\zeta(ns)\) has a pole only for \(n=q\).
Every other \(n/q>0\) is a nonzero finite zeta value by the real-axis
calculation above.  Since \(q\) is prime and \(m\geq2\), its only
divisor at least \(m\) is \(q\), so \(c_m(q)=1/q\).
The residue formula gives the claimed exponent.

A small loop about \(1/q\) therefore has character
\(e^{2\pi i/q}\), of order \(q\).  Infinitely many primes exceed \(m\):
if only finitely many did, all primes would form a finite list, while
one plus their product has a prime divisor outside that list.
Thus the image of \(\chi_m\) is infinite.

On a meromorphic realization over a local chart
\(s-1/q=t^e\), the order is \(e/q\), which must be an integer.
Therefore its ramification index is at least \(q\).
A finite-degree cover has all local indices bounded by its degree.
Arbitrarily large primes contradict such a bound, both on \(\Hp\)
and in any neighborhood of zero containing their reciprocals.
\end{proof}

Multiplying by a single-valued meromorphic factor cannot remove this
obstruction.  Its local order is an integer, so it changes
\(\alpha_m(a)\) by an integer and leaves the character unchanged.
The result concerns this determinant germ and its explicitly constructed
continuation.  It neither identifies a global arithmetic cohomology nor
asserts that the entire imaginary axis is a natural boundary.
The continuation through \(\Hp\), its minimal local ramification,
and its order-comparison maps are determined by the formulas above.

